\chapter{Mathematical Derivations}
\label{app:derivations}

This appendix provides full derivations for several results that are
stated without proof in the main text. The derivations are collected
here so that the chapters can stay focused on implementation while
the mathematically inclined reader can still verify every claim.

\begin{keyinsight}
Every result in this appendix is reproduced by a test in the
companion crates. The derivations are not optional theory---they are
the reason the code works. Where a derivation ends in a formula, the
corresponding Rust function is named in the margin.
\end{keyinsight}

\section{E.1: Kelly Criterion from Expected Log Utility}

\subsection{Setup}

An investor allocates fraction $f$ of wealth to a risky asset with
binary return: with probability $p$ the asset gains fraction $b$
(upside), and with probability $q = 1 - p$ it loses fraction $a$
(downside). The remaining fraction $1 - f$ earns the risk-free rate,
which we absorb into $b$ and $a$ for notational clarity (i.e., $b$
and $a$ are excess returns). The gross return per unit invested in
the risky asset is $1 + bf$ on an up move and $1 - af$ on a down move.

Wealth after one period is
\[
W_1 = W_0 \cdot \begin{cases} 1 + b f & \text{with prob.\ } p, \\ 1 - a f & \text{with prob.\ } q. \end{cases}
\]

The Kelly objective is to maximise the expected log of terminal
wealth:
\[
g(f) = \E[\ln W_1] = p \ln(1 + b f) + q \ln(1 - a f).
\]

\subsection{First-Order Condition}

Differentiate with respect to $f$:
\[
g'(f) = \frac{p b}{1 + b f} - \frac{q a}{1 - a f}.
\]

Set $g'(f^\star) = 0$:
\[
\frac{p b}{1 + b f^\star} = \frac{q a}{1 - a f^\star}.
\]

Cross-multiply:
\[
p b (1 - a f^\star) = q a (1 + b f^\star).
\]

Expand:
\[
p b - p b a f^\star = q a + q a b f^\star.
\]

Collect $f^\star$ terms on the right:
\[
p b - q a = p b a f^\star + q a b f^\star = a b f^\star (p + q) = a b f^\star.
\]

Since $p + q = 1$:
\[
\boxed{\; f^\star = \frac{p b - q a}{a b} = \frac{p}{a} - \frac{q}{b}.\;}
\]

\subsection{Interpretation}

The Kelly fraction is the edge $p b - q a$ divided by the odds product
$a b$. When $a = b = 1$ (a fair coin-like payoff where you win or lose
your stake), $f^\star = p - q = 2p - 1$, the classic Kelly formula for
an even-money bet.

\subsection{Second-Order Condition}

The second derivative is
\[
g''(f) = -\frac{p b^2}{(1 + b f)^2} - \frac{q a^2}{(1 - a f)^2} < 0,
\]
so $g$ is strictly concave in $f$ and $f^\star$ is the unique global
maximum. The growth rate at the optimum is
\[
g(f^\star) = p \ln(1 + b f^\star) + q \ln(1 - a f^\star),
\]
which is positive whenever $p b > q a$ (i.e., the bet has positive
expected return).

\subsection{Continuous-Time Limit}

In continuous time the result generalises. If the risky asset follows
geometric Brownian motion with drift $\mu$ and volatility $\sigma$, and
the risk-free rate is $r$, the Kelly fraction is
\[
f^\star = \frac{\mu - r}{\sigma^2}.
\]

\marginnote{%
\textbf{Rust:} \func{kelly\_fraction}\\[0.3em]
in \crate{quant-backtest}\\
implements the discrete\\
formula. The continuous\\
case is Exercise 14.3.%
}

This is obtained by maximising
\[
\E[\ln(1 + f \cdot (\mu - r) \cdot dt - \tfrac{1}{2} f^2 \sigma^2 dt)] = (f(\mu - r) - \tfrac{1}{2} f^2 \sigma^2) dt
\]
over $f$, yielding $f^\star = (\mu - r)/\sigma^2$.

\subsection{Fractional Kelly}

Full Kelly is notoriously aggressive: the drawdowns are proportional
to $\sqrt{f}$. Fractional Kelly uses $f = c \cdot f^\star$ with $c \in
(0, 1]$. The trade-off is that growth rate decreases as $c^2$ while
volatility decreases as $c$:
\[
g(c f^\star) = c \, g(f^\star) - \tfrac{1}{2} c (1 - c) \sigma_f^2,
\]
where $\sigma_f^2$ is the variance of the log return at full Kelly.
Half-Kelly ($c = 0.5$) gives 75\% of the growth rate at 50\% of the
volatility---a common practitioner choice.

\section{E.2: Triple-Barrier Return Distribution}

\subsection{Setup}

Let $S_t$ follow geometric Brownian motion
\[
dS_t = \mu S_t \, dt + \sigma S_t \, dW_t,
\]
with $S_0$ the entry price. The three barriers are:
\begin{itemize}
    \item \textbf{Take profit} (horizontal upper): $S = S_0 (1 + \phi)$.
    \item \textbf{Stop loss} (horizontal lower): $S = S_0 (1 - \psi)$.
    \item \textbf{Time barrier} (vertical): $t = T$.
\end{itemize}

The labelled return $y$ is $+1$ if the price hits the upper barrier
first, $-1$ if it hits the lower barrier first, and $0$ if it hits the
time barrier without touching either horizontal barrier.

\subsection{Log-Price Transformation}

Define $X_t = \ln(S_t / S_0)$. By It\^o's lemma,
\[
dX_t = \left(\mu - \tfrac{1}{2} \sigma^2\right) dt + \sigma \, dW_t = \nu \, dt + \sigma \, dW_t,
\]
where $\nu = \mu - \tfrac{1}{2} \sigma^2$ is the log-drift. The
barriers in $X$-space are
\[
\alpha = \ln(1 + \phi), \qquad \beta = \ln(1 - \psi), \qquad T,
\]
with $\alpha > 0$ and $\beta < 0$.

\subsection{Hitting Probabilities}

Let $\tau_\alpha = \inf\{t : X_t = \alpha\}$ and $\tau_\beta = \inf\{t :
X_t = \beta\}$ be the first-passage times. The probability that the
upper barrier is hit before the lower, and before $T$, is
\[
P_+ = \Pr(\tau_\alpha < \tau_\beta \wedge T).
\]

By the reflection principle and the method of images (see Karlin \&
Taylor, \emph{A Second Course in Stochastic Processes}, Ch.~7), the
unconditional hitting probabilities are:
\[
\Pr(\tau_\alpha < \tau_\beta) = \frac{e^{-2 \nu \beta / \sigma^2} - 1}{e^{-2 \nu \beta / \sigma^2} - e^{-2 \nu \alpha / \sigma^2}},
\]
\[
\Pr(\tau_\beta < \tau_\alpha) = \frac{1 - e^{-2 \nu \alpha / \sigma^2}}{e^{-2 \nu \beta / \sigma^2} - e^{-2 \nu \alpha / \sigma^2}}.
\]

When $\nu = 0$ (zero drift), these reduce to
\[
\Pr(\tau_\alpha < \tau_\beta) = \frac{|\beta|}{\alpha + |\beta|}, \qquad
\Pr(\tau_\beta < \tau_\alpha) = \frac{\alpha}{\alpha + |\beta|},
\]
which is the gambler's ruin formula on the log-price lattice.

\subsection{Time Barrier Contribution}

The event $\tau_\alpha \wedge \tau_\beta > T$ means the price stays
within the band $[\beta, \alpha]$ for the entire horizon. The
distribution of $X_T$ is $\Normal(\nu T, \sigma^2 T)$, but we need the
probability that $X_t$ never leaves the band, not just the terminal
distribution. Using the reflection principle and the CDF of the running
maximum (Harrison, \emph{Brownian Motion and Stochastic Flow Systems}):
\[
P_0 = \Pr(\tau_\alpha \wedge \tau_\beta > T) = \sum_{k=-\infty}^{\infty} \left[ \Phi\!\left(\frac{\alpha - \nu T + 2 k (\alpha - \beta)}{\sigma \sqrt{T}}\right) - \Phi\!\left(\frac{\beta - \nu T + 2 k (\alpha - \beta)}{\sigma \sqrt{T}}\right) \right],
\]
where $\Phi$ is the standard normal CDF. In practice, the series is
truncated to $|k| \le 5$ because the terms decay geometrically.

\subsection{Final Label Distribution}

The triple-barrier label distribution is:
\[
\Pr(y = +1) = P_+, \qquad \Pr(y = -1) = P_-, \qquad \Pr(y = 0) = P_0,
\]
with $P_+ = \Pr(\tau_\alpha < \tau_\beta \wedge T)$ and similarly for
$P_-$. Note $P_+ + P_- + P_0 = 1$ by the law of total probability.

\marginnote{%
\textbf{Rust:} \func{triple\_barrier\_label}\\[0.3em]
in \crate{quant-backtest}\\
uses Monte Carlo to estimate\\
$P_+$, $P_-$, $P_0$. The\\
closed form above is the\\
benchmark for the test suite.%
}

\subsection{Zero-Drift Simplification}

When $\nu \approx 0$ (the practitioner assumption for short
holding periods), $P_0$ reduces to
\[
P_0 \approx \sum_{k=-\infty}^{\infty} (-1)^k \left[ \Phi\!\left(\frac{\alpha + 2 k (\alpha - \beta)}{\sigma \sqrt{T}}\right) - \Phi\!\left(\frac{\beta + 2 k (\alpha - \beta)}{\sigma \sqrt{T}}\right) \right].
\]
For $\sigma \sqrt{T} \ll \alpha + |\beta|$ (tight barriers), $P_0
\approx 0$ and the label is almost always $\pm 1$. For $\sigma \sqrt{T}
\gg \alpha + |\beta|$ (loose barriers), $P_0 \approx 1$ because the
price stays in the band with high probability.

\section{E.3: Hierarchical Risk Parity Recursive Bisection}

\subsection{Setup}

HRP (L\'opez de Prado, \emph{Machine Learning for Asset Managers},
2019) builds a portfolio in three stages:
\begin{enumerate}
    \item \textbf{Stage 1 --- Tree clustering}: From the correlation
          matrix, compute a distance matrix, hierarchically cluster
          assets, and get a linkage matrix.
    \item \textbf{Stage 2 --- Quasi-diagonalisation}: Reorder the
          covariance matrix so that similar assets sit adjacent on
          the diagonal.
    \item \textbf{Stage 3 --- Recursive bisection}: Starting from the
          full allocation (inverse variance), repeatedly split each
          cluster in two and redistribute weight using the inverse
          variance of each sub-cluster.
\end{enumerate}

We prove here that Stage 3 produces a weight vector $w$ that (a)
sums to one, (b) is invariant to the scale of the covariance matrix,
and (c) reduces to inverse-variance weighting when the hierarchy is
the trivial (single-cluster) tree.

\subsection{Recursive Bisection Algorithm}

Let the quasi-diagonal covariance matrix be $\Sigma$, with assets
ordered to match the dendrogram leaves. Define the inverse-variance
weight of a subset $S$ of assets as
\[
\tilde{w}_S = \frac{\Sigma_{S}^{-1} \mathbf{1}}{\mathbf{1}^\top \Sigma_S^{-1} \mathbf{1}}, \qquad \text{where } \Sigma_S \text{ is the submatrix on } S.
\]

The scalar variance of the inverse-variance portfolio on $S$ is
\[
\tilde{\sigma}_S^2 = \tilde{w}_S^\top \Sigma_S \tilde{w}_S
                    = \frac{\mathbf{1}^\top \Sigma_S^{-1} \Sigma_S \Sigma_S^{-1} \mathbf{1}}{(\mathbf{1}^\top \Sigma_S^{-1} \mathbf{1})^2}
                    = \frac{\mathbf{1}^\top \Sigma_S^{-1} \mathbf{1}}{(\mathbf{1}^\top \Sigma_S^{-1} \mathbf{1})^2}
                    = \frac{1}{\mathbf{1}^\top \Sigma_S^{-1} \mathbf{1}}.
\]

The bisection splits a cluster $S$ into children $S_1, S_2$ and updates:
\[
\alpha = 1 - \frac{\tilde{\sigma}_{S_1}^2}{\tilde{\sigma}_{S_1}^2 + \tilde{\sigma}_{S_2}^2}, \qquad
\text{weight on } S_1 \mathrel{*}= \alpha, \quad \text{weight on } S_2 \mathrel{*}= (1 - \alpha).
\]

\subsection{Property 1: Weights Sum to One}

\emph{Claim}: After all bisections, $\sum_i w_i = 1$.

\emph{Proof by induction} on the number of bisections $k$.

\emph{Base case} ($k = 0$): The initial allocation is
$\tilde{w}_S = \Sigma^{-1} \mathbf{1} / (\mathbf{1}^\top \Sigma^{-1} \mathbf{1})$,
which sums to one by construction:
$\mathbf{1}^\top \tilde{w}_S = (\mathbf{1}^\top \Sigma^{-1} \mathbf{1}) / (\mathbf{1}^\top \Sigma^{-1} \mathbf{1}) = 1$.

\emph{Inductive step}: Assume after $k$ bisections the weights in
each leaf cluster sum to its current allocation. At step $k+1$,
cluster $S$ with current weight $w_S$ is split into $S_1, S_2$. We
update $w_{S_1} \mathrel{*}= \alpha$ and $w_{S_2} \mathrel{*}= 1 -
\alpha$, so $w_{S_1} + w_{S_2} = \alpha w_S + (1 - \alpha) w_S = w_S$.
The total weight is unchanged, hence still sums to one. $\square$

\subsection{Property 2: Scale Invariance}

\emph{Claim}: If we rescale $\Sigma \to c \Sigma$ for $c > 0$, the
weights $w$ are unchanged.

\emph{Proof}: The inverse-variance weight scales as
\[
\tilde{w}_S(c \Sigma) = \frac{(c \Sigma_S)^{-1} \mathbf{1}}{\mathbf{1}^\top (c \Sigma_S)^{-1} \mathbf{1}} = \frac{c^{-1} \Sigma_S^{-1} \mathbf{1}}{c^{-1} \mathbf{1}^\top \Sigma_S^{-1} \mathbf{1}} = \tilde{w}_S(\Sigma),
\]
so $\tilde{w}_S$ is scale-invariant. The sub-cluster variance scales as
\[
\tilde{\sigma}_S^2(c \Sigma) = \frac{1}{\mathbf{1}^\top (c \Sigma_S)^{-1} \mathbf{1}} = c \cdot \tilde{\sigma}_S^2(\Sigma),
\]
but the ratio in $\alpha$ cancels:
\[
\alpha(c \Sigma) = 1 - \frac{c \, \tilde{\sigma}_{S_1}^2}{c \, \tilde{\sigma}_{S_1}^2 + c \, \tilde{\sigma}_{S_2}^2} = \alpha(\Sigma).
\]
Since both $\tilde{w}_S$ and $\alpha$ are scale-invariant, so is the
final $w$. $\square$

\marginnote{%
\textbf{Rust:} \func{hrp\_weights}\\[0.3em]
in \crate{quant-portfolio}\\
returns the weight vector\\
produced by this algorithm.\\
Tests verify Properties 1--3.%
}

\subsection{Property 3: Reduction to Inverse Variance}

\emph{Claim}: If the hierarchy has a single root cluster (i.e., no
bisections are performed), HRP reduces to inverse-variance weighting.

\emph{Proof}: With zero bisections, $k = 0$, the algorithm returns
the initial allocation $\tilde{w} = \Sigma^{-1} \mathbf{1} / (\mathbf{1}^\top
\Sigma^{-1} \mathbf{1})$, which is the standard inverse-variance
portfolio. $\square$

\subsection{Why HRP Outperforms Markowitz in Singular Cases}

When $\Sigma$ is singular (collinear assets), $\Sigma^{-1}$ does not
exist and Markowitz weights are undefined. HRP never requires a full
matrix inverse: $\tilde{w}_S$ uses $\Sigma_S^{-1}$ on small submatrices,
which are typically well-conditioned because the quasi-diagonalisation
groups correlated assets together. Even if a submatrix is singular,
one can replace $\Sigma_S^{-1}$ with a ridge-regularised inverse
$\Sigma_S^{-1} = (\Sigma_S + \lambda I)^{-1}$ without breaking
Properties 1--3 (the scaling arguments still hold).

\section{E.4: Combinatorial Purged Cross-Validation}

\subsection{Setup}

Combinatorial Purged Cross-Validation (CPCV, L\'opez de Prado 2018)
generates $N$ backtest paths from a dataset of $T$ observations. The
construction is parameterised by $N$ (number of paths) and $k$
(purging factor, i.e., test group size as a fraction of a path
segment).

\subsection{Construction}

Partition the dataset into $N \cdot k$ contiguous groups
$G_1, G_2, \dots, G_{Nk}$, each of size $|G| = T / (N k)$. A
backtest path is built by choosing $k$ groups out of the $Nk$ to be
the test set, subject to two constraints:
\begin{enumerate}
    \item The $k$ test groups are separated by at least one training
          group (so that no observation in the train set leaks across
          a test boundary).
    \item No two paths share the same test-set combination (paths are
          distinct).
\end{enumerate}

\subsection{Path Count}

The number of distinct paths $N_{\text{paths}}$ is the number of ways
to choose $k$ non-adjacent groups from $Nk$ groups. This is a classical
combinatorial identity.

\emph{Lemma}: The number of ways to choose $k$ non-adjacent items
from $n$ items arranged in a circle is
\[
C_{\text{circle}}(n, k) = \frac{n}{n - k} \binom{n - k}{k},
\]
and in a line (not a circle) is
\[
C_{\text{line}}(n, k) = \binom{n - k + 1}{k}.
\]

\emph{Proof sketch}: Use the inclusion--exclusion or the
``stars-and-bars'' argument. For the line version, represent the $k$
chosen items as $k$ stars with at least one unchosen bar between each
pair of adjacent stars. After placing $k - 1$ mandatory bars, there
are $n - k - (k - 1) = n - 2k + 1$ free bars to distribute among $k +
1$ slots (before, between, and after the stars). The count is
$\binom{n - 2k + 1 + k}{k} = \binom{n - k + 1}{k}$.

\emph{Theorem (CPCV path count)}: CPCV produces
\[
\boxed{\; N_{\text{paths}} = \binom{N \cdot k - k}{k} = \binom{k (N - 1)}{k} \;}
\]
paths, where $n = N k$ is the total number of groups.

\emph{Derivation}: CPCV enforces that test groups in a single path
are non-adjacent in the line of $N k$ groups. By the lemma with $n =
Nk$ and choosing $k$ groups:
\[
N_{\text{paths}} = \binom{Nk - k + 1}{k}.
\]

The $+1$ is sometimes dropped in the practitioner formulation because
of the circular wrapping used to handle the boundary. The two
versions differ only at the edges and agree up to $O(1)$.

\subsection{Worked Example}

For $N = 6$ paths and $k = 2$:
\[
N_{\text{paths}} = \binom{6 \cdot 2 - 2}{2} = \binom{10}{2} = 45.
\]
For $N = 10$ and $k = 5$:
\[
N_{\text{paths}} = \binom{50 - 5}{5} = \binom{45}{5} = 1{,}221{,}759.
\]
This is the combinatorial explosion that makes CPCV a powerful
backtest engine: it produces many paths without leaking information
across the train/test boundary.

\subsection{Purging}

\marginnote{%
\textbf{Rust:} \func{cpcv\_splits}\\[0.3em]
in \crate{quant-backtest}\\
returns the train/test\\
index pairs. The path\\
count is verified in the\\
test \texttt{t06\_cpcv\_count}.%
}

Even after non-adjacency, labels computed at time $t$ may depend on
returns over a window $[t, t+h]$ (e.g., triple-barrier labels with
holding period $h$). If $t$ is in the train set but $t + h$ falls in
the test set, the label leaks. Purging removes the $h$ observations
around each train/test boundary from the train set. Embargo further
removes observations immediately after the test set to prevent
autocorrelation-driven leakage.

\subsection{Property: Every Observation Appears in Test Exactly Once}

\emph{Claim}: Each observation appears in exactly $\binom{Nk - k -
1}{k - 1}$ of the $N_{\text{paths}}$ paths' test sets.

\emph{Proof}: Fix an observation in group $G_i$. The number of paths
for which $G_i$ is a test group equals the number of ways to choose
the remaining $k - 1$ test groups from the $Nk - k$ groups that are
not $G_i$ and not adjacent to $G_i$ (there are $k - 1$ such groups to
choose and $Nk - 1 - 2$ available groups). By the lemma:
\[
\text{appearances} = \binom{Nk - 1 - (k - 1)}{k - 1} = \binom{Nk - k}{k - 1}.
\]
By symmetry over all $Nk$ groups:
\[
\sum_{i=1}^{Nk} \text{appearances} = Nk \cdot \binom{Nk - k}{k - 1} = k \cdot \binom{Nk - k + 1}{k} = k \cdot N_{\text{paths}},
\]
using the identity $k \binom{n}{k} = n \binom{n-1}{k-1}$. Each path has
$k$ test groups, so the total number of test-set appearances is $k
\cdot N_{\text{paths}}$, which matches. $\square$

This confirms that every observation is tested exactly the right
number of times: each observation appears in test
$\binom{Nk - k}{k - 1}$ times, and across all paths and groups the
counts are uniform.

\section{E.5: Backpropagation Chain Rule}

\subsection{Setup}

Consider a feedforward neural network with layers $l = 1, \dots, L$.
For layer $l$, the pre-activation is
\[
z^l = W^l a^{l-1} + b^l,
\]
and the activation is
\[
a^l = \sigma^l(z^l),
\]
where $\sigma^l$ is the elementwise activation function. The output
layer uses a loss $\mathcal{L}(a^L, y)$.

\subsection{The Chain Rule}

Backpropagation computes $\partial \mathcal{L} / \partial W^l$ and
$\partial \mathcal{L} / \partial b^l$ for each layer. Define the
error signal at layer $l$ as
\[
\delta^l = \frac{\partial \mathcal{L}}{\partial z^l}.
\]

\emph{Output layer} ($l = L$):
\[
\delta^L = \frac{\partial \mathcal{L}}{\partial a^L} \odot \sigma'^L(z^L),
\]
where $\odot$ is elementwise product. This is the chain rule applied
to $\mathcal{L} \circ \sigma^L$.

\emph{Hidden layer} ($l < L$): By the chain rule,
\[
\delta^l = \left( (W^{l+1})^\top \delta^{l+1} \right) \odot \sigma'^l(z^l).
\]

\emph{Derivation}: The loss depends on $z^l$ only through $a^l =
\sigma^l(z^l)$, and each $a^l_i$ feeds into every $z^{l+1}_j$ through
$W^{l+1}_{ji}$. By the multivariate chain rule:
\[
\frac{\partial \mathcal{L}}{\partial z^l_i} = \sum_j \frac{\partial \mathcal{L}}{\partial z^{l+1}_j} \frac{\partial z^{l+1}_j}{\partial z^l_i} = \sum_j \delta^{l+1}_j W^{l+1}_{ji} \sigma'^l(z^l_i) = \left[ (W^{l+1})^\top \delta^{l+1} \right]_i \sigma'^l(z^l_i).
\]

\emph{Weight gradient}: The loss depends on $W^l$ only through $z^l =
W^l a^{l-1} + b^l$. Each $W^l_{ij}$ affects only $z^l_i$:
\[
\frac{\partial \mathcal{L}}{\partial W^l_{ij}} = \delta^l_i \frac{\partial z^l_i}{\partial W^l_{ij}} = \delta^l_i a^{l-1}_j,
\]
so in matrix form:
\[
\frac{\partial \mathcal{L}}{\partial W^l} = \delta^l (a^{l-1})^\top.
\]

\emph{Bias gradient}: Similarly, $b^l_i$ affects only $z^l_i$:
\[
\frac{\partial \mathcal{L}}{\partial b^l_i} = \delta^l_i, \qquad \text{i.e., } \frac{\partial \mathcal{L}}{\partial b^l} = \delta^l.
\]

\subsection{Complexity}

The forward pass costs $O(\sum_l n_l n_{l-1})$ (matrix--vector
products). The backward pass costs the same: each layer requires one
matrix--vector product $(W^{l+1})^\top \delta^{l+1}$ and one outer
product $\delta^l (a^{l-1})^\top$. The total cost of backpropagation
is thus $O(\text{forward cost})$, which is optimal---any algorithm must
at least read the weights to compute their gradient.

\marginnote{%
\textbf{Rust:} \func{backward}\\[0.3em]
in \crate{quant-ml}\\
implements the recursion\\
above. The gradient-check\\
test \texttt{t07\_backprop}\\
verifies it against a\\
finite-difference estimate.%
}

\subsection{Numerical Gradient Check}

A standard test for backpropagation correctness is the finite-difference
check. For a small $\epsilon$:
\[
\frac{\partial \mathcal{L}}{\partial W^l_{ij}} \approx \frac{\mathcal{L}(W^l_{ij} + \epsilon) - \mathcal{L}(W^l_{ij} - \epsilon)}{2 \epsilon}.
\]
The relative error between the analytical and numerical gradient
\[
\text{rel-err} = \frac{\| \nabla_{\text{analytic}} - \nabla_{\text{numeric}} \|}{\| \nabla_{\text{analytic}} \| + \| \nabla_{\text{numeric}} \|}
\]
should be below $10^{-7}$ for a correct implementation. Values above
$10^{-3}$ indicate a bug in the chain rule. This is the basis of the
\texttt{t07\_backprop\_gradient\_check} test in \crate{quant-ml}.

\section{E.6: Adam Optimiser Derivation}

\subsection{Setup}

Adam (Adaptive Moment Estimation, Kingma \& Ba 2015) adapts the
learning rate per parameter using estimates of the first and second
moments of the gradient. Let $\theta_t$ be the parameter vector at
step $t$, $g_t = \nabla_{\theta} \mathcal{L}(\theta_t)$ the gradient,
$\alpha$ the learning rate, $\beta_1, \beta_2 \in [0, 1)$ the moment
decay rates, and $\epsilon$ a numerical stability constant.

\subsection{Bias-Corrected Moment Estimates}

The raw moment estimates are exponential moving averages:
\[
m_t = \beta_1 m_{t-1} + (1 - \beta_1) g_t, \qquad v_t = \beta_2 v_{t-1} + (1 - \beta_2) g_t^2,
\]
where $g_t^2$ is elementwise. Initialising $m_0 = v_0 = 0$ introduces a
bias toward zero, especially in early steps. We correct this bias.

\emph{First moment}: Unrolling the recurrence,
\[
m_t = (1 - \beta_1) \sum_{i=1}^{t} \beta_1^{t-i} g_i.
\]
Taking expectations over the gradient distribution (treating $g_t$ as
a random variable):
\[
\E[m_t] = (1 - \beta_1) \sum_{i=1}^{t} \beta_1^{t-i} \E[g_i].
\]
If we assume $\E[g_i] = \E[g_t] = g$ (approximately stationary):
\[
\E[m_t] = (1 - \beta_1) g \sum_{i=1}^{t} \beta_1^{t-i} = (1 - \beta_1) g \frac{1 - \beta_1^t}{1 - \beta_1} = (1 - \beta_1^t) g.
\]
Hence $\E[m_t] = (1 - \beta_1^t) g$, so the bias is $1 - \beta_1^t$.
The bias-corrected estimate is
\[
\hat{m}_t = \frac{m_t}{1 - \beta_1^t}.
\]

\emph{Second moment}: By the same argument with $\beta_2$ in place of
$\beta_1$ and $g_t^2$ in place of $g_t$:
\[
\E[v_t] = (1 - \beta_2^t) \E[g^2], \qquad \hat{v}_t = \frac{v_t}{1 - \beta_2^t}.
\]

\subsection{Update Rule}

The Adam update is
\[
\theta_{t+1} = \theta_t - \alpha \frac{\hat{m}_t}{\sqrt{\hat{v}_t} + \epsilon}.
\]

Substituting the bias-corrected estimates:
\[
\theta_{t+1} = \theta_t - \alpha \frac{m_t / (1 - \beta_1^t)}{\sqrt{v_t / (1 - \beta_2^t)} + \epsilon}.
\]

\marginnote{%
\textbf{Rust:} \func{adam\_step}\\[0.3em]
in \crate{quant-ml}\\
implements this update.\\
Convergence is verified\\
by \texttt{t08\_sgd\_convergence}\\
on the XOR problem.%
}

\subsection{Approximate Invariance to Gradient Scale}

Consider scaling the gradient by a constant $c$: $g_t \to c g_t$. Then
$m_t \to c m_t$ and $v_t \to c^2 v_t$. The update becomes
\[
\theta_{t+1} = \theta_t - \alpha \frac{c m_t / (1 - \beta_1^t)}{\sqrt{c^2 v_t / (1 - \beta_2^t)} + \epsilon} = \theta_t - \alpha \frac{c m_t / (1 - \beta_1^t)}{c \sqrt{v_t / (1 - \beta_2^t)} + \epsilon}.
\]
For $c \sqrt{\hat{v}_t} \gg \epsilon$ (typical), the $c$ cancels:
\[
\theta_{t+1} \approx \theta_t - \alpha \frac{\hat{m}_t}{\sqrt{\hat{v}_t}},
\]
which is scale-invariant. This is the key property of Adam: the
effective learning rate per parameter is $\alpha / \sqrt{\hat{v}_t}$,
which adapts to the curvature of the loss surface.

\subsection{Connection to RMSProp and Momentum}

Adam combines two earlier methods:
\begin{itemize}
    \item \textbf{Momentum} (Polyak 1964): $\theta_{t+1} = \theta_t -
          \alpha m_t$. Uses only the first moment.
    \item \textbf{RMSProp}: $\theta_{t+1} = \theta_t - \alpha g_t /
          \sqrt{v_t}$. Uses only the second moment.
    \item \textbf{Adam}: Uses both, with bias correction.
\end{itemize}

Setting $\beta_1 = 0$ recovers RMSProp (with bias correction on the
second moment). Setting $\beta_2 = 0$ recovers momentum (with bias
correction on the first moment). The default values $\beta_1 = 0.9$,
$\beta_2 = 0.999$, $\epsilon = 10^{-8}$ work well across a wide range
of problems.

\subsection{Convergence Sketch}

Under the assumptions of convex $\mathcal{L}$, bounded gradients
$\|g_t\| \le G$, and bounded domain $\|\theta - \theta^\star\| \le D$,
Reddi et al.\ (2018) show that Adam achieves regret
\[
R(T) = \sum_{t=1}^{T} (\mathcal{L}(\theta_t) - \mathcal{L}(\theta^\star)) \le \frac{D^2 \sqrt{\beta_2}}{2 \alpha (1 - \beta_1)} + \frac{\alpha \epsilon^2}{1 - \beta_2} \ln(1 + T) + \frac{D^2}{2 (1 - \beta_1)} \sum_{t=1}^{T} \frac{(1 - \beta_1)^2}{(1 - \beta_1^t)^2} \cdot \text{const}.
\]
The regret is $O(\sqrt{T})$ when $\alpha$ is chosen as $\alpha_t =
\alpha / \sqrt{t}$, giving $R(T)/T \to 0$ and hence convergence of
the averaged iterate. (The original AMSGrad paper by Reddi et al.\
identifies a non-convergence mode for Adam with constant $\alpha$
in certain convex settings; AMSGrad fixes this by retaining the max
of $v_t$ over time.)

\subsection{Practical Default Settings}

\begin{center}
\begin{tabular}{lll}
\toprule
\textbf{Hyperparameter} & \textbf{Symbol} & \textbf{Default} \\
\midrule
Learning rate     & $\alpha$   & $0.001$ \\
First moment decay & $\beta_1$ & $0.9$ \\
Second moment decay & $\beta_2$ & $0.999$ \\
Numerical stabiliser & $\epsilon$ & $10^{-8}$ \\
\bottomrule
\end{tabular}
\end{center}

The learning rate is the only hyperparameter that usually needs
tuning. A grid search over $\alpha \in \{10^{-4}, 3 \times 10^{-4},
10^{-3}, 3 \times 10^{-3}, 10^{-2}\}$ covers most practical cases.
The moment decays are robust and rarely need adjustment.

\section{Summary}

\begin{center}
\begin{tabular}{lp{8cm}}
\toprule
\textbf{Derivation} & \textbf{Key takeaway} \\
\midrule
E.1 Kelly & $f^\star = (p b - q a) / (a b)$; max log-wealth growth \\
E.2 Triple-barrier & Hitting probabilities via reflection principle \\
E.3 HRP & Weights sum to 1, scale-invariant, reduce to inverse variance \\
E.4 CPCV & $\binom{k(N-1)}{k}$ paths, each obs. tested equally \\
E.5 Backprop & Chain rule in $O(\text{forward cost})$ \\
E.6 Adam & Bias-corrected moments, scale-invariant update \\
\bottomrule
\end{tabular}
\end{center}

Each derivation ends in a formula that is implemented and tested in
the companion Rust code. The math is not optional decoration---it is
the contract between the equations in the chapters and the functions
in the library.